\section{Метод граничных элементов}
\label{sec:boundary-element-method}

\subsection{Основная идея}

\textbf{Метод граничных элементов (МГЭ)} --- численный метод краевых задач, основанный на переходе от дифференциальной постановки к граничному интегральному уравнению. В классической постановке неизвестные дискретизируются на границе $\Gamma=\partial\Omega$, поэтому размерность сетки уменьшается на единицу:
\[
3D\rightarrow2D,
\qquad
2D\rightarrow1D.
\]

Особенно естественен МГЭ для линейных задач, для оператора которых известно фундаментальное решение: теория потенциала, линейная упругость, акустика и др.

\subsection{Фундаментальное решение и формула Грина}

Пусть для определённости рассматривается оператор Лапласа и выбирается фундаментальное решение
\[
-\Delta G(x,\xi)=\delta(x-\xi).
\]
Тогда
\[
G(x,\xi)=
\begin{cases}
-\dfrac{1}{2\pi}\ln|x-\xi|,& d=2,\\[2mm]
\dfrac{1}{4\pi|x-\xi|},& d=3.
\end{cases}
\]
Знак формул меняется при другой конвенции для оператора, поэтому важно использовать одну согласованную систему знаков.

Для гармонической функции $\Delta u=0$ в $\Omega$ вторая формула Грина приводит к интегральному представлению. Если $q=\partial u/\partial n$, то для внутренней точки $\xi\in\Omega$
\[
\boxed{
u(\xi)=
\int_\Gamma G(x,\xi)q(x)\,d\Gamma
-
\int_\Gamma u(x)\frac{\partial G(x,\xi)}{\partial n_x}\,d\Gamma.}
\]

При предельном переходе точки $\xi$ на границу получается \textbf{граничное интегральное уравнение}
\[
\boxed{
c(\xi)u(\xi)
+
\int_\Gamma u(x)\frac{\partial G(x,\xi)}{\partial n_x}\,d\Gamma
=
\int_\Gamma G(x,\xi)q(x)\,d\Gamma,
\qquad \xi\in\Gamma.}
\]
Для гладкой границы
\[
c(\xi)=\frac12;
\]
в угловой точке коэффициент определяется локальным углом.

Если имеется объёмный источник, например $-\Delta u=f$, в интегральном представлении появляется интеграл по $\Omega$. Поэтому полностью граничное преимущество классического МГЭ непосредственно сохраняется прежде всего для однородных уравнений; объёмный член приходится вычислять или дополнительно преобразовывать.

\subsection{Дискретизация границы}

Граница разбивается на элементы
\[
\Gamma=\bigcup_{e=1}^{N}\Gamma_e.
\]
В двумерной задаче это обычно отрезки, в трёхмерной --- треугольники или четырёхугольники. На каждом элементе аппроксимируются $u$ и $q$:
\[
u(x)\approx\sum_jN_j(x)u_j,
\qquad
q(x)\approx\sum_jN_j(x)q_j.
\]
Простейший вариант --- кусочно-постоянные граничные элементы.

После подстановки и применения, например, метода коллокаций формируется система
\[
\boxed{Hu=Gq+b},
\]
где матрицы $H$ и $G$ состоят из интегралов ядер $\partial G/\partial n$ и $G$, а $b$ содержит вклад объёмных источников, если он есть.

На участках границы известны условия Дирихле $u=\bar u$, Неймана $q=\bar q$ или Робена. После перестановки известных и неизвестных величин получается обычная СЛАУ
\[
Ax=f.
\]
Найдя неизвестные на границе, значение во внутренней точке восстанавливают по интегральному представлению.

\subsection{Сингулярные интегралы}

Ядро фундаментального решения сингулярно при $x\to\xi$. Поэтому диагональные и соседние интегралы нельзя бездумно вычислять обычной квадратурой. Используются аналитическое интегрирование, специальные квадратуры, регуляризация и главное значение интеграла. В локальном члене сингулярного интеграла возникает коэффициент $c(\xi)$.

Корректная обработка особенностей --- одна из центральных технических задач реализации МГЭ.

\subsection{Преимущества и ограничения}

\textbf{Преимущества:}
\begin{itemize}
\item дискретизируется только граница;
\item фундаментальное решение уже учитывает свойства дифференциального оператора;
\item удобно описывать внешние и неограниченные области;
\item граничные напряжения/потоки в ряде постановок определяются с высокой точностью.
\end{itemize}

\textbf{Ограничения:}
\begin{itemize}
\item требуется фундаментальное решение;
\item объёмные источники, неоднородности и сильные нелинейности усложняют чисто граничную постановку;
\item необходимо обрабатывать сингулярные интегралы;
\item матрицы классического МГЭ обычно плотные, в отличие от разреженных матриц МКЭ.
\end{itemize}

\subsection{Сравнение с МКЭ}

МКЭ дискретизирует всю область и строится из слабой/вариационной постановки. МГЭ использует фундаментальное решение и граничное интегральное уравнение, дискретизируя границу. Поэтому МГЭ особенно привлекателен для линейных однородных задач и бесконечных областей, а МКЭ более универсален для неоднородных, нелинейных и сложных многоматериальных задач.

\subsection{Что важно сказать на экзамене}

На экзамене достаточно уверенно провести цепочку
\[
\boxed{
Lu=f
\rightarrow
\text{фундаментальное решение}
\rightarrow
\text{граничное интегральное уравнение}
\rightarrow
\text{граничные элементы}
\rightarrow
Ax=f}.
\]
Обязательно упомянуть коэффициент $c=1/2$ на гладкой границе, сингулярность ядер и плотность матрицы.

\newpage
